\begin{Theorem}\label{theorem:JR}
On a connected smooth projective variety $M$, take a holomorphic Cartan geometry $H\to E \to M$ modelled on an infinitesimally algebraic complex homogeneous space~\pr{X,G}.
Suppose that the induced $G$-bundle $G \to E_G \to M$ also admits a flat holomorphic connection.
Then $M$ belongs to a tower of holomorphic fibrations $M \to M' \to S$ where~$M'$ has no rational curves.
The Cartan geometry is lifted from~$M'$.
The map $M' \to S$ is a holomorphic fibration of abelian varieties.
On every fiber, the Cartan connection is a holomorphic connection on the pullback $G$-bundle with solvable holonomy.
The base $S$ has ample canonical bundle.
\end{Theorem}

\begin{proof}
Replace $M$ by $M'$ without loss of generality, so assume that $M$ contains no rational curves.
Replace $M$ by a finite unramified covering space, so that we can replace $G$ by a finite index subgroup,
and therefore assume that $G$ has a semisimple Levi quotient over $\LieG$
\[
1 \to G_{\rm s} \to G \to G_{\rm ss} \to 1.
\]
The flat holomorphic connection induces a flat holomorphic connection on the quotient $G_{\rm ss} \to E_{G_{\rm ss}} :=
 E/G_{\rm s}\to M$.
Let $L \subset G_{\rm ss}$ be the complex algebraic Zariski closure of the holonomy group of that connection.
This connection induces a flat holomorphic connection on a principal subbundle $L\to E_L \to M$,
$E_L \subset E_{G_{\rm ss}}$.
It therefore also induces a flat holomorphic connection on the quotient bundle $L_{\rm ss} \to E_{L_{\rm ss}}
 = E_L/\sqrt{L} \to M$, where $1 \to \sqrt{L} \to L \longrightarrow
L_{\rm ss} \to 1$ is the semisimple Levi quotient \cite[Theorem 4.3]{Hochschild:1981}.
Denote by~${\rho \colon \fundamentalGroup{M} \to L_{\rm ss}}$ the induced representation.

After perhaps replacing $M$ by a finite unramified covering space, the representation $\rho \colon
\fundamentalGroup{M}\allowbreak\to L_{\rm ss}$ lifts from a big representation $\fundamentalGroup{S} \to L_{\rm ss}$,
where $S := \Shaf[\rho]{M}$ and $S$ is a projective variety of
general type \cite[Theorem 5 (and remarks following)]{Zuo:1999}.
The Shafarevich fibration~${\rationalMap{M}{S}}$ is defined and proper on a Zariski dense open subset of $M$
\cite[Theorem 4.1]{Kollar:1993}.
The general fiber has finite unramified covering on which the holonomy of $E_{L_{\rm ss}}\to M$ is trivial, i.e., the holonomy of $E_L \to M$ lies in $\sqrt{L}$, i.e., the holonomy of the $E_G$-connection lies in the preimage of $\sqrt{L}$ in $G$, i.e., an extension of $\sqrt{L}$ by $G_{\rm s}$.
By Lemma~\ref{lemma:abelian.pancake}, the general fiber admits a~finite unramified covering by an abelian variety, with abelian holonomy image and $\cb{M}$ is trivial on the general fiber.

The general fiber of $\rationalMap{M}{S}$ has a finite unramified covering by an abelian variety, so the Iitaka conjecture is verified for $\rationalMap{M}{S}$: $\kappa_M \ge \kappa_S = \dim{S}$.
Since $\kappa_M \ge \dim{S}\ge 0$, we know that $M$ has an Iitaka fibration.
The general fiber $M_{s}$ of $\rationalMap{M}{S}$ admits a finite unramified covering by an abelian variety and $\cb{M}$ is trivial on $M_{s}$, so $M_{s}$ lies in a fiber of the Iitaka fibration, giving a rational map $\rationalMap{S}{\Ii{M}}$ making a commutative diagram
\[
\begin{tikzcd}
M \arrow[<->]{r}\arrow[->, dashed,densely dashed]{d} & M \arrow[->, dashed,densely dashed]{d} \\
S \arrow[->,dashed,densely dashed]{r} & \Ii{M}.
\end{tikzcd}
\]
The rational map $\rationalMap{S}{\Ii{M}}$ is dominant, as the top row of the commutative diagram is dominant.
By dominance, $\dim{S} \ge \dim{\Ii{M}} = \kappa_M \ge \kappa_S = \dim S$.
Consequently $\rationalMap{S}{\Ii{M}}$ is a~birational map identifying $\rationalMap{M}{S}$ with the Iitaka fibration
\cite[Theorem 10.6]{Iitaka:1982}, i.e., the Shafarevich fibration is a model of the Iitaka fibration.

Since the generic fiber of the Iitaka fibration has canonical bundle spanned by global sections, so does the total
space \cite[Theorem 4.4]{Lai:2011}, i.e., the canonical bundle of $M$ is spanned by global sections.
Hence the Iitaka fibration is holomorphic, i.e., we can arrange that $M \to S$ is holomorphic.

The canonical bundle of $M$ restricts to a nef line bundle on every fiber of $M \to S$, since~$M$ has canonical bundle spanned by global sections.
Some power of the canonical bundle $K_M$ is pulled back to $M$ from $M \to S$ \cite[Theorem 2.1.26]{Lazarsfeld:2004}.
Hence that power is trivial on the fibers, and so $\acb{M}$ is also nef on every fiber of $M \to S$.
Therefore, all fibers of $M \to S$ have the same dimension \cite[Theorem 2]{Kawamata:1991}.

Take any resolution of singularities $S_0 \to S$ and pullback $M$ to get a dominant morphism~${M_0 \to S_0}$
birational to $M \to S$ with smooth $M_0$ and $S_0$.
After replacing $M$ with a~finite unramified covering space, there is some resolution of singularities $S_0 \longrightarrow
S$ over which~${M \to S}$ is birational to a holomorphic fibration (indeed, a smooth morphism, which is also a $C^{\infty}$ fiber bundle \cite[Definition 5.7]{Kollar:1993}) with fibers abelian varieties, with smooth total space, over a smooth variety of general type, equipped with a holomorphic
section \cite[Theorem~6.3]{Kollar:1993}:
\[
\begin{tikzcd}
\dot{M} \arrow[->]{d} \arrow[<->, dashed,densely dashed]{r} & M_0 \arrow[->]{r}\arrow[->]{d} & M \arrow[->]{d}\\
\dot{S} \arrow[<->, dashed,densely dashed]{r} & S_0 \arrow[->]{r} & S.
\end{tikzcd}
\]
Rational maps to $M$ extend to holomorphic maps, since $M$ contains no rational curves.
Extend the rational map $\rationalMap{\dot{M}}{M}$ to a holomorphic map $\dot{M} \to M$.
Up to birational isomorphism, there is a section $\dot{S} \longrightarrow
 \dot{M}$ \cite[Theorem 6.3]{Kollar:1993}, and so a rational section $\rationalMap{S}{M}$.
Extend the rational section to a holomorphic map $S \to M$, and hence a holomorphic section.
This section maps curves in $S$ to isomorphic curves in $M$; but $M$ contains no rational curves, and so $S$ contains no rational curves.
Hence meromorphic maps to $S$ extend to be holomorphic.
Holomorphically extend $\rationalMap{\dot{S}}{S}$ and $\rationalMap{M}{S}$:
\[
\begin{tikzcd}
\dot{M} \arrow[->,shift left]{r}\arrow[->]{d} & M \arrow[->]{d} \arrow[->, dashed,densely dashed,shift left]{l} \\
\dot{S} \arrow[->,shift left]{r} & S\arrow[->, dashed,densely dashed,shift left]{l}.
\end{tikzcd}
\]

By Zariski's main theorem (see \cite[Corollary 11.4]{Hartshorne:1977}), the map $\dot{M} \to M$ has connected fibers, as does $\dot{S} \to S$.
Hence $\dot{M} \to \dot{S}$ is also an Iitaka fibration.
Each fiber of $\dot{M}\to \dot{S}$ is an abelian variety mapping holomorphically to a fiber of $M \to S$.

On the generic fiber $\dot{M}_{\dot{s}}$, this map $\dot{M}_{\dot{s}}\to M_{s}$ is birational with exceptional locus covered in rational curves.
But neither fiber contains rational curves.
Therefore, the generic fiber $\dot{M}_{\dot{s}}$ of~$\dot{M}\to \dot{S}$ is mapped isomorphically to a fiber
of $M\to S$.

Suitable powers of the canonical bundles of $\dot{M}$ and $M$ are spanned by global sections, pulled back from sections of line bundles on $\dot{S}$ and $S$ respectively.
All sections on $M$ and on $\dot{M}$ are pulled back from $S$ and $\dot{S}$ respectively.
So the vanishing locus of any such section is a union of fibers.

The holomorphic map $\dot{M} \to M$ is an isomorphism except on the ramification divisor, an effective divisor given by the vanishing of the determinant of the derivative of $\dot{M} \to M$.
Applying the derivative of $\dot{M} \to M$ to sections of the canonical bundle on $M$, we get sections of the canonical bundle on $\dot{M}$, with divisor the sum of the pullback canonical divisor and the ramification divisor.
But the sections vanish on unions of fibers.
Hence the ramification divisor of $\dot{M}\to M$ is the preimage of some effective divisor on $\dot{S}$.

Away from the image in $S$ of the support of that divisor in $\dot{S}$, our section of $M \to S$ strikes a smooth fiber in a single smooth point, since $\dot{M} \to M$ is biholomorphic on those fibers.
All fibers have intersection number 1 with the section, and so all fibers are smooth near the image of the section, and the image of the section is smooth.
The variety $S$ is smooth, since the image of the section is isomorphic to $S$.

Since the fibers of $M \to S$ are all of the same dimension, at any point of any fiber at which the reduction of the fiber is smooth, we can pick a transverse complex submanifold, a~local holomorphic section of $M \to S$.
All nearby fibers have intersection number 1 with the local section, and so with our fiber: the generic point of every fiber of $M
\to S$ is reduced.

If some fiber $M_{s}$ of $M \longrightarrow
S$ is reducible, its preimage in $\dot{M}$ is reducible.
As the fibers of~${\dot{M} \longrightarrow
M}$ are connected, the preimage in $\dot{M}$ of $M_{s}$ is a connected reducible subvariety.
Since the fibers of $\dot{M} \to \dot{S}$ are irreducible, the preimage of $M_{s}$ in $\dot{M}$ is a union of fibers over a~connected reducible subvariety of $\dot{S}$.
This subvariety of $\dot{S}$ maps to a point in $S$.
Because the subvariety is connected, all of the fibers over the subvariety map to the same component of~$M_{s}$.
But $\dot{M} \to M$ is surjective, since it is holomorphic and birational.
Therefore, all fibers of~$M \to S$ are irreducible.

The fibers $\dot{M}_{\dot{s}}$ of $\dot{M} \to \dot{S}$ all have the same images in homology inside $M$, all nonzero as they are algebraic cycles.
So the derivative of the holomorphic map $\dot{M}_{\dot{s}} \to M$ drops rank on a ramification divisor.
All holomorphic sections of a suitable positive power of the canonical bundle of $M$ are pulled back from sections of a suitable line bundle on $S$.
Find a section of that suitable line bundle, nonzero at $s$, and pullback to get a section of $\cb{M}$ nowhere zero near $M_{s}$.
Take a local holomorphic section of the canonical bundle of $S$ defined and nonzero at $s$, and divide the two sections to get a holomorphic section of the relative canonical bundle of $M \to S$ defined near $M_{s}$.
Plugging the derivative of $\dot{M}_{\dot{s}} \to M$ into this section gives a section of the canonical bundle of the abelian variety $\dot{M}_{\dot{s}}$, not vanishing somewhere and therefore not vanishing anywhere as the canonical bundle of an abelian variety is trivial.
Therefore, $\dot{M}_{\dot{s}} \to M$ is a~local biholomorphism on every fiber, taking each fiber $\dot{M}_{\dot{s}}$ to a fiber $M_{s}$.
The map $\dot{M} \to M$ is birational, so the map $\dot{M}_{\dot{s}} \to M_{s}$ on each fiber is a biholomorphism to its image.

Take a local holomorphic section of $M \longrightarrow
 S$ transverse to the image of a fiber $\dot{M}_{\dot{s}}$.
The section strikes the generic nearby fiber of $M \to S$ in a single smooth point, since $\dot{M} \longrightarrow
M$ is biholomorphic on generic fibers.
Generic fibers have intersection number 1 with the section, and so all fibers have intersection 1 with the section, so are smooth and reduced near the image of the section.
Hence all fibers of $M \longrightarrow
S$ are everywhere smooth and reduced and irreducible, and so the maps $\dot{M}_{\dot{s}} \to M_{s}$ are all isomorphisms.

Since $S$ has big canonical bundle, its Iitaka fibration $S \longrightarrow
 \Ii{S}$ is birational, with exceptional fibers covered by rational curves \cite{Kawamata:1991}.
But $S$ contains no rational curves, so there are no exceptional fibers.
Since rational maps to $S$ extend to become holomorphic, the Iitaka fibration~${S \longrightarrow
 \Ii{S}}$ has a holomorphic inverse.
The exceptional locus of the canonical morphism is then also covered by rational curves, so is empty, so $S$ has ample canonical bundle.
\end{proof}

\section{Finite holonomy}

\begin{Example}\label{example:finite.holonomy}
Suppose that \pr{X',G} is a complex homogeneous space and that $X'$ is a complex torus $X' = \C^{n}/\Lambda$.
Take a finite index subgroup $\Lambda_0 \subset \Lambda$.
Let $M' := \C^{n}/\Lambda_0 \to X' = \C^{n}/\Lambda$ be the obvious projection, and pullback the geometry to $M'$, i.e., let $E \longrightarrow
 M'$ be the pullback $H$-bundle
\[
\begin{tikzcd}
E \arrow{r}\arrow{d} & G \arrow{d} \\
M' \arrow{r} & X'.
\end{tikzcd}
\]
The map $M' \longrightarrow
 X'$ is the developing map of a flat holomorphic Cartan geometry modelled on~$(X', G)$: the pullback of the model geometry.
The holonomy morphism has finite image in $G$.
Suppose that $X \to X'$ is a $G$-equivariant holomorphic map of complex homogeneous $G$-spaces, a complex
homogeneous bundle over the complex torus, say $X = G/H$ and $X' = G/H'$, and suppose that $H/H'$ is a flag variety.
Then the quotient $M := E/H$ is the lift of $E \to M'$ to a~flat holomorphic $(X, G)$-geometry.
\end{Example}

\begin{Theorem}\label{theorem:finite.holonomy}
Take a smooth projective variety $M$ and a flat holomorphic Cartan geometry on $M$ modelled on a complex homogeneous space \pr{X,G}.
If the holonomy morphism of the geometry has finite image then, after perhaps replacing $M$ by a finite unramified covering space, the geometry is constructed as in the above Example~{\rm \ref{example:finite.holonomy}}.
\end{Theorem}

\begin{proof}
We can drop to assume that $M$ is not uniruled.
We can assume that the model $X$ is connected.
Since the Cartan connection is flat, we can take a developing map and associated holonomy morphism.
Lift to a finite unramified covering space to arrange trivial holonomy, so the geometry is the pullback via a developing map
$\widehat{f} \colon M \to X$.
Pulling back, ${TM \cong \widehat{f}^* TX}$.
Since the model $X$ is homogeneous, $TX$ is spanned by global sections and so the anticanonical bundle~$\acb{X}$ of $X$ is spanned by global sections.
The canonical bundle of $M$ is nef.
But~${\cb{M} = \widehat{f}^* \cb{X}}$.
Therefore, the canonical bundle of $M$ is trivial.
So the developing map has image inside an anticanonical fiber.
Since the developing map is a local biholomorphism, the anticanonical fibers are open sets in $X$.
By compactness of $M$, the developing map is a finite unramified covering map to a component of $X$.
By Lemma~\ref{lemma:is.a.hoop}, the developing map is a hoop.
So $X$ is an abelian variety~${X = \C^{n}/\Lambda}$.
The group $G$ can be any complex Lie group acting on $X$, while the effectivization is a finite extension of $X$ by some linear transformations of $\C^{n}$ preserving $\Lambda$.
The developing map is an isogeny of abelian varieties $M \to X$.
\end{proof}

\section{Which subvarieties develop to the model}
We sum up our results into a theorem from which we easily obtain Theorem~\ref{theorem:main}.
\begin{Theorem}\label{theorem:big.1}
Take a non-uniruled smooth projective variety $M$ and a flat holomorphic Cartan geometry on $M$ modelled on a complex homogeneous space \pr{X,G}, where $G$ is a complex linear algebraic group.
After perhaps replacing $M$ by a finite unramified covering space, $M$ is a~holomorphic fibration $M \to \grave{M}
 \longrightarrow
 M'$ of abelian varieties over a smooth complex projective variety $\grave{M}$, refining the fibration of Theorem~{\rm\ref{theorem:JR}}, so that a finite map $Z \longrightarrow
 M$ from a~connected projective variety develops to the model $($after perhaps replacing $Z$ by a finite unramified covering$)$ just when the image of $Z \longrightarrow
 M$ lies in a fiber of $M \longrightarrow \grave{M}$.
Every fiber of~${M \to \grave{M}}$ has a finite unramified covering space which develops to a hoop in the model.
The map $M \to \grave{M}$ has a holomorphic section.
\end{Theorem}

\begin{proof}
If the holonomy is finite, apply Theorem~\ref{theorem:finite.holonomy}.
Suppose that the holonomy is infinite.
Let $\Sigma \subset \fundamentalGroup{M}$ be the kernel of the holonomy morphism $\fundamentalGroup{M} \to G$.
Any subvariety $Z \subset M$ whose fundamental group lies in $\Sigma$ has a developing map $\widehat{Z} \longrightarrow
 X$ from its normalization.
The developing map identifies the ambient tangent bundles $ TX|_{\widehat{Z}} = TM|_{\widehat{Z}}$.
Since $M$ contains no rational curves, its canonical bundle is nef, while the tangent bundle of $X$ is spanned by global sections, so both canonical bundles restrict to be trivial along $\widehat{Z}$.
By Lemma~\ref{lemma:dev.large}, $Z \to M$ lies in a leaf of the ant foliation, while $\widehat{Z} \to X$ lies in a fiber of the ant fibration.
Moreover, $Z$ has large fundamental group.

Every smooth fiber $M_{s}$ of $\Shaf[\Sigma]{} \colon \rationalMap{M}{S := \Shaf[\Sigma]{M}}$ has a finite unramified covering which develops to the model, by definition of $\Shaf[\Sigma]{}$.
By Lemma~\ref{lemma:abelian.pancake}, $M_{s}$ has a finite covering by an abelian variety.
By Corollary~\ref{corollary:hoops}, some finite covering of $M_{s}$ maps to $X$ by a hoop.
Every fiber~$M_{s}$ develops to $X$, so if $Z$ lies inside a smooth fiber $M_{s}$, then $Z$ develops to $X$.
Let~${M \to \grave{M}}$ be~${\Shaf[\Sigma]{}\colon M \to \Shaf[\Sigma]{M}}$.
The map $\Shaf[\Sigma]{} \colon M \to \Shaf[\Sigma]{M}$ comes with a rational section,
say $\rationalMap[s]{\grave{M}}{M}$ \cite[Theorem~6.3]{Kollar:1993}.
Meromorphic maps to $M$ extend to a holomorphic maps, since $M$ contains no rational curves.
\end{proof}
